In HE3DB,
the LWE ciphertexts are filtered results 
with a small dimension ($n=1024$) and modulus ($q$ is bit).
The straight application of [25] results in packed RLWE ciphertext 
also having small parameters, and we cannot perform further aggregation 
on such small-parameter RLWE ciphertext.

First of all, 
we would like to thank all the reviewers for their valuable comments. 
In what follows, we summarize our responses to the questions,
and plan to revise our manuscript accordingly.
*** 
# Review #118A 
**Question A-1**: Analyze complexity between HE3DB and [73].
**Answer A-1**:  
We find out that our method requires $1+(2n/5)$ bootstrapping operations
for the MSB extraction of an $n$-bit encrypted plaintext,
while [73] requires $1+(2n/3)$ (in the fresh-noise settings, it
becomes $1+2n/4$). Therefore, our method enjoys a better slope than 
the method in [73], which translates to up to 1.6x faster for 32-bit 
comparison for DB.
Besides, we also have different decomposition parameter settings, and
we use FFT instead of NTT, which can all result in performance differences.
We will revise the manuscript to include the above discussions and
correct the mistakes.   

*** 
# Review #118B
**Question B-1**: The cost of achieving malicious security over FHE?
**Answer B-1**:
Achieving malicious security in an FHE-based approach indeed
requires additional cryptographic techniques 
such as zero-knowledge proofs or OPRF (e.g., [a]). 
However, current methods do result in relatively high computational 
costs, and can be slower than TEE or MPC-based approaches. 
Despite these challenges, 
we believe that providing malicious security for homomorphic encryption 
is an important and promising direction of research.

[a] Hao Chen et al., "Labeled PSI from Fully Homomorphic Encryption with 
Malicious Security", CCS, 2018.

**Question B-2**: It seems access pattern is leaked in the current setting. Can MPC or TEE be better?
**Answer B-2**: 
We would first like to kindly point out that, similar to HEDA [91], 
HE$^3$DB performs a linear scan of the entire database for each query,
thereby ensuring the privacy of the query access pattern.
We agree that there are also MPC or TEE-based approaches which can
hide access patterns [96, 16]. 
However, they rely on additional interactions or tools, 
which can result in increased communication costs or slower performance.
We will add more clarifications
on these approaches in Section 1.2 of the main manuscript.

**Question B-3**: :Can you comment on why HEDA provides 80-bit security and you provide 128-bit?
**Answer B-3**:
We apologize that this is actually a mistake. HEDA does provide 128-bit
security, and we will correct this mistake in the revised manuscript.

*** 
# Review #118C

**Question C-1**: Alg 1. looks overcomplicated. Can't we simply ... to obtain ct_4?
**Answer C-1**: 
As discussed in Section 1 of [31] (Chillotti et al.) 
and Figure 2 of [75] (Lu et al.),
to correctly bootstrap a TFHE ciphertext, 
we have to ensure that the MSB of the encrypted plaintext is zero.
Therefore, before applying PBS on $ct_1$, 
it is necessary to discard the MSB of the plaintext
encrypted in $ct_1$, which is basically why we need to first produce
ct_4.
We will correct the main manuscript to include the above explanations.

**Question C-2**: Isn't possible ... to directly 
evaluate the function? The proposed method ... faster and more accurate?
**Answer C-2**:
Different from HMSB, the result of a high-precision absolute function 
depends on all bits of the encrypted plaintext.
Consequently, we cannot use the shift-and-subtract strategy in Alg. 2 
to remove the lower bits as noises like HMSB.
For homomorphic min/max and sorting functions, we show that we are faster
than the open-source implementations of [28], and are both faster
and more accurate than [75] (Please kindly refer to Fig. 6 and Fig. 7),
which are two of the most recent state-of-the-art methods.

**Question C-3**: The given LWEstoRLWE algorithm may be slower than
[25].
**Answer C-3**:
While [25] is also an excellent way of converting LWEs to an RLWE,
we find two limitations 
in applying [25] to encrypted databases.
First, it is not known how to change parameters (e.g., $n$ and $q$)
when converting a set of LWE ciphertexts to an RLWE ciphertext using [25].
Second, the runtime of HomMod is independent of the number of LWE ciphertexts
(e.g., $N/2$), while the runtime [25] grows linearly with $N$ 
(therefore O($N$)).
As a result, in the application
of encrypted DB, we find that HomMod is more suitable and
faster than [25] in practice. 
We will add more explanations on
the applicability of [25] in the revised manuscript.


**Comment C-4**: Range of the plaintext?
**Answer C-4**:
We encoded the plaintext as integers in the range $[-2^{k-1}, 2^{k-1})$ and 
adopt the two's complement representation where the MSB of the value indicates 
where the number is positive or negative (the sign).
In two's complement, the MSB of $a-b$ represents the sign bit of the 
result, which is 0 if $a\geq b$,
and 1 for if $a  <  b$. We verified the correctness
of the proposed comparison in the submitted codes.

**Comment C-5**: Most solutions can be ... the parameter. 
**Answer C-5**: We agree that [28] does support 32-bit precision, and
many methods can support unbounded depth using standard FHE bootstrapping. 
We will correct the mistakes and revise the expressions
in Table 3 to better explain the practical benefits of our method 
when applied in DB.